\chapter{Test und Verifikation}

\begin{lernziele}
Nach diesem Kapitel können Sie Anforderungen mit Testfällen verknüpfen und einfache Verifikationsszenarien beschreiben. Sie verstehen den Unterschied zwischen Verifikation und Validierung und können Testfälle für Navigation, Hinderniserkennung, Personenerkennung, Batteriemanagement und Ladestation formulieren.
\end{lernziele}

\section{Einleitung}
Ein Systemmodell ist nur dann wirklich nützlich, wenn es bis zum Nachweis führt. Anforderungen sollen nicht nur schön formuliert werden. Am Ende muss geprüft werden, ob das System sie erfüllt. Genau darum geht es in diesem Kapitel.

\section{Verifikation und Validierung}
Verifikation fragt: Wurde das System richtig gebaut? Das bedeutet: Erfüllt das System die spezifizierten Anforderungen?

Validierung fragt: Wurde das richtige System gebaut? Das bedeutet: Löst das System tatsächlich das Problem der Stakeholder?

Beide Fragen sind wichtig. Ein Roboter kann alle spezifizierten Anforderungen erfüllen und trotzdem für Nutzer unpraktisch sein. Umgekehrt kann ein Prototyp nützlich wirken, aber wichtige Anforderungen nicht nachweisbar erfüllen.

Das klassische V-Modell veranschaulicht diesen Zusammenhang: Jede Spezifikationsebene auf dem absteigenden Ast hat eine entsprechende Verifikations- oder Validierungsebene auf dem aufsteigenden Ast (siehe Abbildung~\ref{fig:kap19-v-modell}).

\begin{figure}[H]
  \centering
  \adjustbox{max width=\linewidth,max totalheight=0.6\textheight}{% =====================================================================
% Abbildung: V-Modell: Verifikation und Validierung
% Passend zu Kapitel 19, Abschnitt "Verifikation und Validierung"
%
% Einbindung im Buch, z. B. in chapters/19_tests_verifikation.tex:
%
%   \begin{figure}[H]
%     \centering
%     % =====================================================================
% Abbildung: V-Modell: Verifikation und Validierung
% Passend zu Kapitel 19, Abschnitt "Verifikation und Validierung"
%
% Einbindung im Buch, z. B. in chapters/19_tests_verifikation.tex:
%
%   \begin{figure}[H]
%     \centering
%     % =====================================================================
% Abbildung: V-Modell: Verifikation und Validierung
% Passend zu Kapitel 19, Abschnitt "Verifikation und Validierung"
%
% Einbindung im Buch, z. B. in chapters/19_tests_verifikation.tex:
%
%   \begin{figure}[H]
%     \centering
%     \input{figures/fig_kap19_v_modell}
%     \caption{Das V-Modell verbindet den absteigenden Spezifikationsast
%       mit dem aufsteigenden Verifikationsast. Gestrichelte Linien zeigen,
%       welche Testebene welche Spezifikationsebene nachweist.}
%     \label{fig:kap19-v-modell}
%   \end{figure}
%
% Benoetigte Pakete (im Buch bereits in main.tex geladen): tikz
% Benoetigte TikZ-Bibliotheken: positioning, arrows.meta, shadows, calc
% =====================================================================

\resizebox{\textwidth}{!}{\input{figures/fig_kap19_v_modell.tikz}}

%     \caption{Das V-Modell verbindet den absteigenden Spezifikationsast
%       mit dem aufsteigenden Verifikationsast. Gestrichelte Linien zeigen,
%       welche Testebene welche Spezifikationsebene nachweist.}
%     \label{fig:kap19-v-modell}
%   \end{figure}
%
% Benoetigte Pakete (im Buch bereits in main.tex geladen): tikz
% Benoetigte TikZ-Bibliotheken: positioning, arrows.meta, shadows, calc
% =====================================================================

\resizebox{\textwidth}{!}{\begin{tikzpicture}[
    every node/.style={font=\small},
    reqbox/.style={
        rectangle, rounded corners=2pt, draw=black!70, thick,
        minimum width=32mm, minimum height=11mm,
        align=center, fill=blue!12, drop shadow
    },
    archbox/.style={
        rectangle, rounded corners=2pt, draw=green!40!black, thick,
        minimum width=32mm, minimum height=11mm,
        align=center, fill=green!15, drop shadow
    },
    implbox/.style={
        rectangle, rounded corners=2pt, draw=orange!60!black, thick,
        minimum width=32mm, minimum height=11mm,
        align=center, fill=orange!15, drop shadow
    },
    testbox/.style={
        rectangle, rounded corners=2pt, draw=black!50, thick,
        minimum width=32mm, minimum height=11mm,
        align=center, fill=black!6, drop shadow
    },
    flow/.style={-{Stealth[length=2.2mm]}, thick, black!60},
    corr/.style={dashed, thick, black!50},
    lbl/.style={font=\tiny, fill=white, inner sep=1pt}
]

% --- linker Ast: Spezifikation (absteigend) ---
% Hinweis: Abstaende so gewaehlt, dass Boxen (Breite 32mm) auf keiner Ebene
% ueberlappen; der horizontale Versatz zur Mittelachse (x=7) waechst pro
% Ebene um 2cm, sodass auch die engste Ebene (comp/unit, Versatz 2cm,
% Abstand 4cm) mehr Abstand hat als die Boxbreite.
\node[reqbox]  (stake) at (-1, 12) {Stakeholder-\\Bedürfnisse};
\node[reqbox]  (sysreq) at (1, 9)  {System-\\anforderungen};
\node[archbox] (arch)  at (3, 6)  {Architektur/\\Design};
\node[archbox] (comp)  at (5, 3)  {Komponenten-\\entwurf};

% --- Spitze ---
\node[implbox] (impl) at (7, 0) {Implementierung};

% --- rechter Ast: Verifikation (aufsteigend) ---
\node[testbox] (unit)  at (9, 3)  {Unit-Test};
\node[testbox] (integ) at (11, 6) {Integrations-\\test};
\node[testbox] (systest) at (13, 9) {Systemtest};
\node[testbox] (valid) at (15, 12) {Abnahme-/\\Validierungstest};

% --- absteigender Ast ---
\draw[flow] (stake)  -- (sysreq);
\draw[flow] (sysreq) -- (arch);
\draw[flow] (arch)   -- (comp);
\draw[flow] (comp)   -- (impl);

% --- aufsteigender Ast ---
\draw[flow] (impl)    -- (unit);
\draw[flow] (unit)    -- (integ);
\draw[flow] (integ)   -- (systest);
\draw[flow] (systest) -- (valid);

% --- horizontale Entsprechungen (gestrichelt, ohne Pfeilspitze) ---
\draw[corr] (stake)  -- node[lbl, midway]{validiert}   (valid);
\draw[corr] (sysreq) -- node[lbl, midway]{verifiziert} (systest);
\draw[corr] (arch)   -- (integ);
\draw[corr] (comp)   -- (unit);

% --- Randbeschriftungen ---
\node[font=\small, rotate=90, black!70] at (-3.4, 6) {Spezifikation};
\node[font=\small, rotate=-90, black!70] at (17.4, 6) {Verifikation};
\node[font=\small, black!70] at (7, -2.2) {Implementierung};

% --- Legende ---
\node[reqbox,  minimum width=6mm, minimum height=4mm, font=\tiny] (leg1) at (1,-5) {};
\node[right=1mm of leg1, font=\tiny, align=left] {Anforderung};
\node[archbox, minimum width=6mm, minimum height=4mm, font=\tiny, right=24mm of leg1] (leg2) {};
\node[right=1mm of leg2, font=\tiny, align=left] {Architektur/Entwurf};
\node[implbox, minimum width=6mm, minimum height=4mm, font=\tiny, right=24mm of leg2] (leg3) {};
\node[right=1mm of leg3, font=\tiny, align=left] {Implementierung};
\node[testbox, minimum width=6mm, minimum height=4mm, font=\tiny, right=24mm of leg3] (leg4) {};
\node[right=1mm of leg4, font=\tiny, align=left] {Testebene};

\end{tikzpicture}
}

%     \caption{Das V-Modell verbindet den absteigenden Spezifikationsast
%       mit dem aufsteigenden Verifikationsast. Gestrichelte Linien zeigen,
%       welche Testebene welche Spezifikationsebene nachweist.}
%     \label{fig:kap19-v-modell}
%   \end{figure}
%
% Benoetigte Pakete (im Buch bereits in main.tex geladen): tikz
% Benoetigte TikZ-Bibliotheken: positioning, arrows.meta, shadows, calc
% =====================================================================

\resizebox{\textwidth}{!}{\begin{tikzpicture}[
    every node/.style={font=\small},
    reqbox/.style={
        rectangle, rounded corners=2pt, draw=black!70, thick,
        minimum width=32mm, minimum height=11mm,
        align=center, fill=blue!12, drop shadow
    },
    archbox/.style={
        rectangle, rounded corners=2pt, draw=green!40!black, thick,
        minimum width=32mm, minimum height=11mm,
        align=center, fill=green!15, drop shadow
    },
    implbox/.style={
        rectangle, rounded corners=2pt, draw=orange!60!black, thick,
        minimum width=32mm, minimum height=11mm,
        align=center, fill=orange!15, drop shadow
    },
    testbox/.style={
        rectangle, rounded corners=2pt, draw=black!50, thick,
        minimum width=32mm, minimum height=11mm,
        align=center, fill=black!6, drop shadow
    },
    flow/.style={-{Stealth[length=2.2mm]}, thick, black!60},
    corr/.style={dashed, thick, black!50},
    lbl/.style={font=\tiny, fill=white, inner sep=1pt}
]

% --- linker Ast: Spezifikation (absteigend) ---
% Hinweis: Abstaende so gewaehlt, dass Boxen (Breite 32mm) auf keiner Ebene
% ueberlappen; der horizontale Versatz zur Mittelachse (x=7) waechst pro
% Ebene um 2cm, sodass auch die engste Ebene (comp/unit, Versatz 2cm,
% Abstand 4cm) mehr Abstand hat als die Boxbreite.
\node[reqbox]  (stake) at (-1, 12) {Stakeholder-\\Bedürfnisse};
\node[reqbox]  (sysreq) at (1, 9)  {System-\\anforderungen};
\node[archbox] (arch)  at (3, 6)  {Architektur/\\Design};
\node[archbox] (comp)  at (5, 3)  {Komponenten-\\entwurf};

% --- Spitze ---
\node[implbox] (impl) at (7, 0) {Implementierung};

% --- rechter Ast: Verifikation (aufsteigend) ---
\node[testbox] (unit)  at (9, 3)  {Unit-Test};
\node[testbox] (integ) at (11, 6) {Integrations-\\test};
\node[testbox] (systest) at (13, 9) {Systemtest};
\node[testbox] (valid) at (15, 12) {Abnahme-/\\Validierungstest};

% --- absteigender Ast ---
\draw[flow] (stake)  -- (sysreq);
\draw[flow] (sysreq) -- (arch);
\draw[flow] (arch)   -- (comp);
\draw[flow] (comp)   -- (impl);

% --- aufsteigender Ast ---
\draw[flow] (impl)    -- (unit);
\draw[flow] (unit)    -- (integ);
\draw[flow] (integ)   -- (systest);
\draw[flow] (systest) -- (valid);

% --- horizontale Entsprechungen (gestrichelt, ohne Pfeilspitze) ---
\draw[corr] (stake)  -- node[lbl, midway]{validiert}   (valid);
\draw[corr] (sysreq) -- node[lbl, midway]{verifiziert} (systest);
\draw[corr] (arch)   -- (integ);
\draw[corr] (comp)   -- (unit);

% --- Randbeschriftungen ---
\node[font=\small, rotate=90, black!70] at (-3.4, 6) {Spezifikation};
\node[font=\small, rotate=-90, black!70] at (17.4, 6) {Verifikation};
\node[font=\small, black!70] at (7, -2.2) {Implementierung};

% --- Legende ---
\node[reqbox,  minimum width=6mm, minimum height=4mm, font=\tiny] (leg1) at (1,-5) {};
\node[right=1mm of leg1, font=\tiny, align=left] {Anforderung};
\node[archbox, minimum width=6mm, minimum height=4mm, font=\tiny, right=24mm of leg1] (leg2) {};
\node[right=1mm of leg2, font=\tiny, align=left] {Architektur/Entwurf};
\node[implbox, minimum width=6mm, minimum height=4mm, font=\tiny, right=24mm of leg2] (leg3) {};
\node[right=1mm of leg3, font=\tiny, align=left] {Implementierung};
\node[testbox, minimum width=6mm, minimum height=4mm, font=\tiny, right=24mm of leg3] (leg4) {};
\node[right=1mm of leg4, font=\tiny, align=left] {Testebene};

\end{tikzpicture}
}
}
  \caption{Das V-Modell verbindet den absteigenden Spezifikationsast mit dem aufsteigenden Verifikationsast.}
  \label{fig:kap19-v-modell}
\end{figure}

\section{Testfälle}
Die folgende Tabelle listet die Testfälle des Roboterprojekts auf. Das Kürzel \enquote{TC} steht dabei für \emph{Test Case} (Testfall), das Kürzel \enquote{REQ} für \emph{Requirement} (Anforderung, siehe Kapitel 6).

\begin{longtable}{p{0.16\textwidth}p{0.34\textwidth}p{0.34\textwidth}}
\toprule
ID & Testfall & Verifizierte Anforderung\\
\midrule
TC-001 & Autonome Navigation zu Zielpunkt & REQ-001\\
TC-002 & Hindernis rechtzeitig erkennen & REQ-002\\
TC-003 & Hindernis sicher behandeln & REQ-003\\
TC-004 & Person im Kamerabild erkennen & REQ-004\\
TC-005 & Niedrigen Akkustand melden & REQ-005\\
TC-006 & Zur Ladestation fahren & REQ-006\\
TC-007 & Betriebszustand melden & REQ-007\\
TC-008 & Sensorausfall sicher behandeln & REQ-008\\
\bottomrule
\end{longtable}

\section{Testbeschreibung}
Ein Testfall sollte enthalten:

\begin{itemize}
    \item ID
    \item Ziel
    \item verifizierte Anforderung
    \item Vorbedingungen
    \item Testaufbau
    \item Ablauf
    \item erwartetes Ergebnis
    \item Bewertungskriterium
    \item benötigte Messdaten
\end{itemize}

\section{Beispiel TC-002}
\textbf{Vorbedingung:} Roboter fährt mit definierter Geschwindigkeit.\\
\textbf{Ablauf:} Hindernis wird in den Fahrweg gestellt.\\
\textbf{Erwartung:} Roboter erkennt Hindernis und stoppt vor Kollision.\\
\textbf{Kriterium:} Mindestabstand 20 cm.

\section{Testarten}
\subsection{Unit-Tests}
Unit-Tests prüfen kleine Softwareeinheiten. Ein Batteriemonitor kann zum Beispiel getestet werden, indem simulierte Spannungswerte eingegeben und erwartete Akkustände geprüft werden.

\subsection{Integrationstests}
Integrationstests prüfen das Zusammenspiel mehrerer Komponenten. Beispiel: \texttt{battery\_monitor\_node} meldet niedrigen Akkustand, \texttt{navigation\_node} plant Rückkehr zur Ladestation.

\subsection{Systemtests}
Systemtests prüfen das Verhalten des gesamten Roboters in definierten Szenarien.

\subsection{Simulationstests}
Simulationstests erlauben viele Tests ohne echte Hardware. Sie ersetzen reale Tests nicht vollständig, sind aber für frühe Entwicklung sehr wertvoll. Kapitel 22 stellt mit Gazebo ein konkretes Simulationswerkzeug für solche Tests vor.

\section{Tests im Modell}
In SysML können Testfälle mit \texttt{verify}-Beziehungen an Anforderungen gehängt werden. Dadurch ist erkennbar, welche Anforderungen noch keine Tests haben. Abbildung~\ref{fig:kap19-verify-beziehung} zeigt dies am Beispiel von TC-006 und REQ-006.

In SysML v2 kann eine solche \texttt{verify}-Beziehung textuell etwa so notiert werden:

\begin{lstlisting}
package Tests {
  import Requirements::*;

  test def TC_006 {
    doc /* Prueft, ob der Roboter bei niedrigem Akkustand
           selbststaendig zur Ladestation faehrt. */
    verify requirement REQ_006;
  }
}
\end{lstlisting}

\begin{figure}[H]
  \centering
  \adjustbox{max width=\linewidth,max totalheight=0.6\textheight}{% =====================================================================
% Abbildung: SysML-verify-Beziehung zwischen Testfall und Anforderung
% Passend zu Kapitel 19, Abschnitt "Tests im Modell"
%
% Einbindung im Buch, z. B. in chapters/19_tests_verifikation.tex:
%
%   \begin{figure}[H]
%     \centering
%     % =====================================================================
% Abbildung: SysML-verify-Beziehung zwischen Testfall und Anforderung
% Passend zu Kapitel 19, Abschnitt "Tests im Modell"
%
% Einbindung im Buch, z. B. in chapters/19_tests_verifikation.tex:
%
%   \begin{figure}[H]
%     \centering
%     % =====================================================================
% Abbildung: SysML-verify-Beziehung zwischen Testfall und Anforderung
% Passend zu Kapitel 19, Abschnitt "Tests im Modell"
%
% Einbindung im Buch, z. B. in chapters/19_tests_verifikation.tex:
%
%   \begin{figure}[H]
%     \centering
%     \input{figures/fig_kap19_verify_beziehung}
%     \caption{Der Testfall TC-006 verifiziert die Anforderung REQ-006
%       über eine \texttt{verify}-Beziehung.}
%     \label{fig:kap19-verify-beziehung}
%   \end{figure}
%
% Benoetigte Pakete (im Buch bereits in main.tex geladen): tikz
% Benoetigte TikZ-Bibliotheken: positioning, arrows.meta, shadows, calc
% =====================================================================

\input{figures/fig_kap19_verify_beziehung.tikz}

%     \caption{Der Testfall TC-006 verifiziert die Anforderung REQ-006
%       über eine \texttt{verify}-Beziehung.}
%     \label{fig:kap19-verify-beziehung}
%   \end{figure}
%
% Benoetigte Pakete (im Buch bereits in main.tex geladen): tikz
% Benoetigte TikZ-Bibliotheken: positioning, arrows.meta, shadows, calc
% =====================================================================

\begin{tikzpicture}[
    every node/.style={font=\small},
    req/.style={
        rectangle, rounded corners=2pt, draw=black!70, thick,
        minimum width=55mm, minimum height=12mm,
        align=center, fill=blue!12, drop shadow
    },
    test/.style={
        rectangle, rounded corners=2pt, draw=black!50, thick,
        minimum width=55mm, minimum height=12mm,
        align=center, fill=black!6, drop shadow
    },
    dependency/.style={
        dashed, thick, black!60,
        -{Triangle[open, length=2.6mm, width=2.2mm]}
    }
]

% --- Knoten ---
\node[req]  (req006) at (0, 3) {\texttt{«requirement»}\\REQ-006: Rückkehr\\zur Ladestation};
\node[test] (tc006)  at (0, -1.5) {\texttt{«test case»}\\TC-006: Zur Ladestation\\fahren};

% --- verify-Beziehung (gestrichelt, offene Pfeilspitze, von Test zu Anforderung) ---
\draw[dependency] (tc006) -- node[midway, fill=white, inner sep=1pt, font=\small\itshape]{«verify»} (req006);

\end{tikzpicture}


%     \caption{Der Testfall TC-006 verifiziert die Anforderung REQ-006
%       über eine \texttt{verify}-Beziehung.}
%     \label{fig:kap19-verify-beziehung}
%   \end{figure}
%
% Benoetigte Pakete (im Buch bereits in main.tex geladen): tikz
% Benoetigte TikZ-Bibliotheken: positioning, arrows.meta, shadows, calc
% =====================================================================

\begin{tikzpicture}[
    every node/.style={font=\small},
    req/.style={
        rectangle, rounded corners=2pt, draw=black!70, thick,
        minimum width=55mm, minimum height=12mm,
        align=center, fill=blue!12, drop shadow
    },
    test/.style={
        rectangle, rounded corners=2pt, draw=black!50, thick,
        minimum width=55mm, minimum height=12mm,
        align=center, fill=black!6, drop shadow
    },
    dependency/.style={
        dashed, thick, black!60,
        -{Triangle[open, length=2.6mm, width=2.2mm]}
    }
]

% --- Knoten ---
\node[req]  (req006) at (0, 3) {\texttt{«requirement»}\\REQ-006: Rückkehr\\zur Ladestation};
\node[test] (tc006)  at (0, -1.5) {\texttt{«test case»}\\TC-006: Zur Ladestation\\fahren};

% --- verify-Beziehung (gestrichelt, offene Pfeilspitze, von Test zu Anforderung) ---
\draw[dependency] (tc006) -- node[midway, fill=white, inner sep=1pt, font=\small\itshape]{«verify»} (req006);

\end{tikzpicture}

}
  \caption{Der Testfall TC-006 verifiziert die Anforderung REQ-006 über eine \texttt{verify}-Beziehung.}
  \label{fig:kap19-verify-beziehung}
\end{figure}

\begin{praxisbox}
Wenn REQ-006 eine Rückkehr zur Ladestation fordert, sollte ein Testfall prüfen, ob der Roboter bei unterschrittenem Akkuschwellwert tatsächlich ein Rückkehrziel setzt und keinen neuen normalen Transportauftrag startet.
\end{praxisbox}

\section{Typische Fehler}
\subsection{Tests zu spät planen}
Wenn Tests erst am Ende entstehen, sind Anforderungen oft unprüfbar.

\subsection{Nur den Normalfall testen}
Fehlerfälle wie Sensorausfall, blockierter Weg und niedriger Akku sind besonders wichtig.

\subsection{Keine klaren Kriterien}
\enquote{Roboter stoppt rechtzeitig} muss in ein messbares Kriterium übersetzt werden.

\section{Testplanung über den Projektverlauf}
Tests sollten nicht erst am Ende des Projekts entstehen. Bereits beim Formulieren einer Anforderung sollte gefragt werden: Wie wird diese Anforderung später geprüft? Dadurch erkennt man früh, ob eine Anforderung zu unklar ist.

Für das Roboterprojekt kann die Testplanung schrittweise wachsen:

\begin{itemize}
    \item Anforderungen erhalten erste Akzeptanzkriterien.
    \item Wichtige Funktionen erhalten Testideen.
    \item Schnittstellen erhalten Integrationstests.
    \item ROS2-Nodes erhalten Unit- oder Komponententests.
    \item Szenarien wie Zielnavigation werden als Systemtests beschrieben.
\end{itemize}

Sinnvoll ist außerdem, bestehende Tests nach jeder Änderung erneut auszuführen (Regressionstest), damit neu eingeführte Fehler frühzeitig auffallen – besonders relevant im Zusammenhang mit dem Änderungsmanagement aus Kapitel 20.

Ein gutes Testmodell macht auch Lücken sichtbar. Wenn eine Anforderung keinen Testfall besitzt, ist das nicht automatisch falsch, aber es ist eine offene Frage. Vielleicht wird sie durch Analyse, Review oder Inspektion nachgewiesen. Test, Analyse, Inspektion und Demonstration gelten dabei als die vier klassischen Verifikationsmethoden. Auch das sollte dokumentiert werden.

\section{Vertiefungsbeispiel: Rückkehrabsicherung prüfen}
\label{sec:bm-tests}
Die Anforderungen BM-01 bis BM-13 aus Abschnitt~\ref{sec:bm-anforderungen} werden durch die folgenden Testentwürfe konkretisiert. Sie erweitern TC-005 und TC-006. Die Tests sind geplant, noch nicht durchgeführt; eine Zuordnung ist kein bestandener Nachweis. Die Kennung T-01 bis T-12 gilt innerhalb dieses Batteriemanagement-Beispiels.

\subsection{Gemeinsamer Testaufbau und Messdaten}
Für die ersten Entscheidungstests werden Akkudaten, Verbrauchsverlauf, Restauftrag, Wegdaten und Stationsstatus kontrolliert vorgegeben, beispielsweise in einer Simulation. Jede Prüfung dokumentiert die verwendeten Reserven, Unsicherheitsgrenzen und Zeitvorgaben. Aufzuzeichnen sind Eingangsdaten mit Zeitstempel und Gültigkeit, Bedarfsschätzungen, Entscheidungsgrund, freigegebene Route, Auftragsstatus sowie Andock- und Ladebestätigung.

Die Prüfung der Entscheidung wird von der Prüfung der Schätzqualität getrennt: Ein korrekter Vergleich vorgegebener Energiewerte beweist nicht, dass die Werte für den realen Roboter stimmen. Wegverfolgung, Energieverbrauch und Andocken benötigen ergänzende Integrations- und Systemtests. Solange die offenen Parameter aus Abschnitt~\ref{sec:bm-anforderungen} fehlen, sind insbesondere Aussagen wie \enquote{rechtzeitig} noch nicht abschließend quantitativ prüfbar.

\subsection{Szenarien und erwartete Ergebnisse}
\begin{description}[style=nextline]
    \item[T-01: Auftrag energetisch gedeckt] Alle Startvoraussetzungen sind erfüllt; Restauftrag, anschließende Rückkehr und Reserve sind gedeckt. Auch die aktuelle Rückkehr ist abgesichert. Erwartung: Start und Fortsetzung werden freigegeben; Schätzwerte und Grund sind nachvollziehbar. Prüft BM-01 bis BM-04 und BM-12.
    \item[T-02: Energie nur für die Rückkehr] Die aktuelle Rückkehr einschließlich Reserve ist gedeckt, der Restauftrag jedoch nicht. Erwartung: Auftrag kontrolliert unterbrechen und Rückkehr beginnen; keine weitere Auftragsfortsetzung. Prüft BM-03, BM-04 und BM-08.
    \item[T-03: Fehlende Karte, aufgezeichneter Rückweg] Karte und Restauftragsschätzung fehlen. Der aufgezeichnete Rückweg ist verfolgbar, befahrbar und energetisch gedeckt. Erwartung: Auch bei hohem Akkustand keine Auftragsfortsetzung; Rückkehr über den aufgezeichneten Weg. Prüft BM-05 bis BM-07.
    \item[T-04: Rückweg nach Auftragsende unbekannt] Der Restauftrag ist bekannt, sein anschließender Rückweg nicht. Vom aktuellen Standort ist eine Rückkehr abgesichert. Erwartung: Auftrag unterbrechen und unmittelbar zurückkehren; unbekannten Bedarf nicht als null behandeln. Prüft BM-03 bis BM-05.
    \item[T-05: Verbrauch oder Rückwegbedarf steigt] Während des Auftrags steigen Verbrauch beziehungsweise Rückwegbedarf. Erwartung: Schätzung bei relevanter Änderung und innerhalb des festgelegten Prüfabstands aktualisieren; Rückkehr vor Verlust der Absicherung einleiten. Danach eine günstigere Schätzung vorgeben: Der Auftrag bleibt während der Rückkehr unterbrochen. Prüft BM-02 bis BM-04, BM-08 und den Ablauf aus Abschnitt~\ref{sec:bm-ablauf}.
    \item[T-06: Wegaufzeichnung fällt aus] Variante A: Eine unabhängig abgesicherte Route bleibt verfügbar. Variante B: Keine solche Route bleibt verfügbar. Erwartung: A führt zur Rückkehr über die gültige Route; B zum kontrollierten Anhalten und zur Unterstützungsanforderung. Lücken dürfen nicht als gültiger Rückweg behandelt werden. Prüft BM-06, BM-07 und BM-10 bis BM-12.
    \item[T-07: Akkudaten ungültig] Es ist keine belastbare Restenergiebewertung mehr möglich. Erwartung: Keine weitere autonome Fahrt freigeben; kontrolliert anhalten und Unterstützung anfordern. Prüft BM-03, BM-04 und BM-10 bis BM-12.
    \item[T-08: Blockade mit tragfähiger Alternative] Der Rückweg wird blockiert. Eine verfolg- und befahrbare Alternative einschließlich Warte-, Planungs- und Andockenergie ist gedeckt. Erwartung: Nur diese geprüfte Alternative freigeben und die Rückkehr darüber fortsetzen. Prüft BM-07 bis BM-09.
    \item[T-09: Blockade ohne tragfähige Alternative] Es gibt keinen abgesicherten alternativen Rückweg oder die Ladestation ist unerreichbar. Erwartung: Kontrolliert anhalten, Auftrag unterbrochen lassen, Unterstützung anfordern und autonome Rückkehr als nicht erfüllt melden. Prüft BM-09 bis BM-12.
    \item[T-10: Ladebeginn fehlt] Die Stationsposition wird erreicht, aber Andocken oder Ladebeginn schlägt fehl. Erwartung: Keine Erfolgsmeldung. Wenn keine abgesicherte Fortsetzung möglich ist, Fehlerstatus und Unterstützung anfordern. Im erfolgreichen Vergleichslauf müssen Andock- und Ladebestätigung beide vorliegen. Prüft BM-10, BM-12 und BM-13.
    \item[T-11: Startvoraussetzung fehlt] Vor Fahrtbeginn fehlen Akkudaten, Verbrauchserfassung oder Rückkehrstrategie; jede Voraussetzung separat ausfallen lassen. Erwartung: Kein Transportauftrag startet; fehlende Voraussetzung wird gemeldet. Prüft BM-01 und BM-12.
    \item[T-12: Entscheidungsgrenze] Gültige Testwerte unmittelbar oberhalb, auf und unterhalb der festgelegten Fortsetzungsgrenze vorgeben. Erwartung: Oberhalb und auf der Grenze darf nur bei erfüllten übrigen Bedingungen fortgesetzt werden; unterhalb wird zurückgekehrt. Die aktuelle Rückkehrgrenze gesondert prüfen: Rechtzeitige Rückkehr hat Vorrang vor Fortsetzung. Prüft BM-03, BM-04 und BM-08.
\end{description}

\subsection{Ein Testfall im Detail: T-03}
\begin{description}[style=nextline]
    \item[Ziel und Vorbedingungen] BM-05 bis BM-07 prüfen. Der Roboter hat einen laufenden Auftrag, gültige Akkudaten und einen unabhängig von der Karte aufgezeichneten, freigegebenen Rückweg. Seine Energie reicht für diesen Weg einschließlich Reserve.
    \item[Testaufbau und Ablauf] In einer kontrollierten Umgebung die Kartendaten und Restauftragsschätzung als nicht verfügbar kennzeichnen; Wegaufzeichnung und ihre Rückkehrbewertung gültig lassen. Anschließend die nächste Auftragsentscheidung auslösen.
    \item[Erwartetes Ergebnis] Der Roboter unterbricht den Auftrag, nennt die fehlende Abschätzbarkeit als Grund und gibt die abgesicherte Rückkehr frei.
    \item[Bewertungskriterium] Nach Verarbeitung der ungültigen Schätzung wird keine weitere Auftragsfortsetzung freigegeben. Die Rückkehr beginnt innerhalb der zuvor festgelegten Reaktionszeit. Kontrolliertes Abbremsen oder eine notwendige Unterbrechungsaktion ist von weiterer Auftragsausführung zu unterscheiden.
    \item[Messdaten] Gültigkeitskennzeichen und Zeitstempel, Energiebedarf und Reserve, ausgewählte Rückwegkennung, Freigaben, Auftragsstatus und Entscheidungsgrund aufzeichnen.
\end{description}

\section{Zusammenfassung}
Tests und Verifikation schließen den Kreis von Anforderungen zur Umsetzung. Ein gutes Modell zeigt nicht nur, was gebaut werden soll, sondern auch, wie die Erfüllung geprüft wird.

\section{Kontrollfragen}
\begin{enumerate}
    \item Was ist Verifikation?
    \item Was ist Validierung?
    \item Welche Informationen enthält ein guter Testfall?
    \item Wozu dient die SysML-Beziehung \texttt{verify}?
    \item Was ist der Unterschied zwischen Unit-Test und Systemtest?
    \item Warum sind Fehlerfälle wichtig?
\end{enumerate}

\section{Übungen}
\begin{uebungbox}
\textbf{Übung 1: Testfälle erstellen}\\
Erstellen Sie Testfälle für alle Anforderungen aus Kapitel 6. Markieren Sie Anforderungen ohne Test.
\end{uebungbox}

\begin{uebungbox}
\textbf{Übung 2: Testfall detaillieren}\\
Beschreiben Sie TC-006 \enquote{Zur Ladestation fahren} mit Vorbedingungen, Ablauf, Erwartung und Bewertungskriterium.
\end{uebungbox}

\section{Literaturhinweise}
Verifikation und Validierung sind zentrale Themen im Systems Engineering. SysML-Literatur \cite{friedenthal2014} zeigt, wie Testfälle mit Anforderungen verbunden werden können.
